\begin{tikzpicture}[font=\small,line cap=round,line join=round]\begin{scope}[scale=.85]\draw[geoTrait](0,0)--(2.6,0)--(1.9,1.4)--(0,1.4)--cycle;\draw[otNavy](.24,0)--(.24,.24)--(0,.24);\draw[otNavy](0,1.16)--(.24,1.16)--(.24,1.4);\node[below,font=\small]at(1.3,-.1){(1)};\begin{scope}[shift={(3.6,0)}]\draw[geoTrait](0,.4)--(2.2,0)--(2.5,1.4)--(.6,1.5)--cycle;\node[below,font=\small]at(1.25,-.1){(2)};\end{scope}\begin{scope}[shift={(7.2,0)}]\draw[geoTrait](0,0)--(2.8,0)--(2.1,1.4)--(.7,1.4)--cycle;\node[below,font=\small]at(1.4,-.1){(3)};\end{scope}\begin{scope}[shift={(10.6,0)}]\draw[geoTrait](0,.2)--(2.0,0)--(1.7,1.5)--(.2,1.2)--cycle;\node[below,font=\small]at(1.0,-.1){(4)};\end{scope}\end{scope}\end{tikzpicture}\par\vspace{1.5mm}{\otRegular\fontsize{8}{10}\selectfont\color{otMuted}Chercher une paire de côtés parallèles : l’écart constant la révèle.\par}
